\documentclass[12pt,a4paper]{article}

\usepackage{fontawesome5}
\providecommand{\faPencilAlt}{\faIcon{pencil-alt}}

\usepackage{styles/couleurs}
\usepackage{styles/mise_en_page}
\usepackage{styles/boites}
\usepackage{styles/maths}
\usepackage{styles/figures}

\usepackage{enumitem}
\usepackage{array}
\usepackage{tabularx}
\usepackage{amsmath}
\usepackage{amssymb}
\usepackage{graphicx}
\usepackage{tikz}
\usepackage{xparse}

\usetikzlibrary{calc,angles,quotes,arrows.meta}

\setlist[itemize]{leftmargin=1.4em,itemsep=0.15em,topsep=0.2em}
\setlist[enumerate]{leftmargin=*,itemsep=0.22em,topsep=0.25em}

\newcommand{\lignepointillee}{\noindent\dotfill\par\vspace{0.45em}}

\newcommand{\lignesreponse}[1]{%
\begin{tcolorbox}[
enhanced,colback=white,colframe=NeutreC!40,boxrule=0.5pt,
arc=1.5mm,left=3mm,right=3mm,top=2mm,bottom=1mm,
before skip=2pt,after skip=2pt]
\foreach \i in {1,...,#1}{\lignepointillee}
\end{tcolorbox}}

\newcommand{\titreseq}[1]{%
\begin{center}
{\sffamily\large\bfseries\color{SavoirC}#1}
\end{center}}

\begin{document}

\titreseq{Activité 6 -- Laboratoire d'erreurs et rédaction}

\section*{Partie A -- Une rédaction attendue}

On veut calculer \(AB\) dans le triangle \(ABC\) rectangle en \(B\),
avec \(AC=9\) cm et \(\widehat A=36^\circ\).

Remets les lignes dans l'ordre.

\begin{itemize}
  \item[\(\square\)] Donc \(AB\approx7{,}3\) cm.
  \item[\(\square\)] \(\cos(\widehat A)=AB\div AC\).
  \item[\(\square\)] \(AB=9\times\cos 36^\circ\).
  \item[\(\square\)] Dans le triangle \(ABC\) rectangle en \(B\), on a :
\end{itemize}

\section*{Partie B -- Quatre copies à analyser}

Pour chaque copie :

\begin{enumerate}
  \item repère la première erreur ;
  \item explique-la ;
  \item propose une correction.
\end{enumerate}

\subsection*{Copie 1 -- Mauvais(e) ...........}

Dans le triangle \(DEF\) rectangle en \(E\), on connaît
\(DF=10\) cm et \(\widehat D=40^\circ\).
On cherche \(DE\).

\[
\cos 40^\circ=EF\div DF.
\]

\[
DE=10\times\cos 40^\circ.
\]

\subsection*{Copie 2 -- Mauvais(e) ...........}

Dans le triangle \(GHI\) rectangle en \(H\), on connaît
\(GH=6\) cm et \(\widehat G=52^\circ\).
On cherche l'hypoténuse \(GI\).

\[
GI=6\times\cos 52^\circ\approx3{,}7\text{ cm}.
\]

\subsection*{Copie 3 -- Mauvais(e) ...........}

Dans une figure composée, un élève écrit directement :

\[
\cos(\widehat A)=BC\div AC,
\]

sans indiquer dans quel triangle il travaille ni vérifier qu'il est rectangle.

\subsection*{Copie 4 -- La calculatrice}

Pour calculer \(8\times\cos 60^\circ\), la calculatrice affiche environ \(7{,}62\).

L'élève conclut :

\[
AB\approx7{,}6\text{ cm}.
\]

\lignesreponse{6}

\section*{Partie C -- Grille de contrôle}

Voici les six questions à se poser avant de valider une solution.

\begin{tcolorbox}[
  colback=SavoirC!5,
  colframe=SavoirC!60,
  boxrule=0.7pt,
  arc=2mm
]
\begin{enumerate}
  \item Ai-je nommé le triangle \textbf{rectangle} ?
  \item Ai-je précisé où se trouve l'angle \textbf{droit} ?
  \item Ai-je correctement repéré l'\textbf{hypoténuse} ?
  \item Ai-je choisi le bon côté \textbf{adjacent à l'angle étudié} ?
  \item Ai-je identifié si la longueur cherchée est
        \textbf{le côté adjacent ou l'hypoténuse} ?
  \item Ai-je choisi la bonne opération :
        \textbf{multiplier ou diviser par le cosinus} ?
  \item Mon résultat est-il cohérent avec le fait que l'hypoténuse
        est le côté le plus \textbf{long} ?
\end{enumerate}
\end{tcolorbox}

\section*{Partie D -- Rédiger entièrement}

Dans le triangle \(JKL\) rectangle en \(K\),
on connaît \(JK=4{,}8\) cm et \(\widehat J=27^\circ\).

Calcule l'hypoténuse \(JL\), arrondie au dixième de centimètre,
en respectant la rédaction attendue.

\lignesreponse{10}

\end{document}
